\newpage{}
\section{Conclusiones y Trabajos Futuros}
\label{sec:conclusiones}

El desarrollo del presente TFG ha culminado con la construcción y despliegue exitoso de \textbf{MoodNest}, una plataforma integral que da respuesta directa a la problemática planteada al inicio de esta memoria: la falta de herramientas de autoconocimiento emocional basadas en datos objetivos que eviten la imprecisión y los sesgos de la memoria retrospectiva.

A continuación, se presenta un balance crítico del proyecto, analizando los resultados obtenidos desde múltiples perspectivas, así como las posibles líneas de evolución de la plataforma.

\subsection{Cumplimiento de Objetivos}

Se confirma de manera rotunda la consecución del objetivo principal del proyecto, ya que MoodNest ofrece una solución que equilibra una interfaz sencilla para el usuario con un potente motor analítico. 

La decisión de implementar un catálogo de etiquetas gestionado en su totalidad por el usuario ha resultado ser un gran acierto para la flexibilidad del sistema. Además, permitir que el usuario ponga una nota específica a cada actividad del día (y no solo una nota general a toda la jornada) marca la diferencia. Esto hace que el sistema calcule automáticamente las estadísticas y dibuje gráficos muy precisos. De este modo, la aplicación le enseña directamente al usuario qué actividades influyen de forma positiva o negativa en su rutina, evitando que dependa de lo que buenamente recuerde o de su propia intuición.
\subsection{Éxito Arquitectónico y Tecnológico}

Desde el punto de vista de la ingeniería de software, el proyecto ha logrado integrar con éxito un \textit{stack} tecnológico moderno, robusto y altamente demandado en el entorno empresarial actual: React.js para la SPA reactiva, Java Spring Boot para la API REST, y MongoDB para la persistencia NoSQL de alta flexibilidad. Asimismo, se ha garantizado la seguridad perimetral de la información confidencial mediante la implementación de una arquitectura \textit{stateless} basada en tokens JWT.

Cabe destacar el éxito en el despliegue del sistema mediante la contenerización con \textbf{Docker}. A pesar de las dificultades inherentes a la orquestación de infraestructuras complejas —tales como los retos técnicos encontrados al configurar las redes internas para conectar el contenedor del servidor Spring Boot con el clúster de la base de datos MongoDB—, el ecosistema fue configurado, depurado y puesto en marcha con éxito, asegurando un entorno predecible y escalable.

\subsection{Éxito Metodológico y Balance Personal}

A nivel de gestión, la adopción de una metodología ágil puramente iterativa ha sido determinante. El uso estricto del sistema de \textit{Issues}, \textit{Milestones} y tableros Kanban a través de GitHub ha permitido mantener un control absoluto sobre el alcance del proyecto. A pesar de contar con un margen temporal ajustado, la priorización continua de tareas garantizó que el TFG llegara a su finalización a tiempo y cumpliendo con los estándares de calidad definidos.

En el plano personal, este TFG ha supuesto un salto cualitativo significativo en mi formación. El principal objetivo formativo era enfrentarme al desarrollo \textit{Full-Stack} completo utilizando herramientas orientadas a entornos empresariales reales, las cuales no siempre se exploran en profundidad durante el curso académico estándar. Sin embargo, gracias a la sólida base de ingeniería de software adquirida durante el Grado, la adaptación a estas nuevas tecnologías ha sido rápida y natural, resultando en un producto completamente funcional del que me siento profundamente orgulloso y que consolida mi perfil técnico de cara al mercado laboral.

\subsection{Líneas de Trabajo Futuro}

MoodNest ha sido concebido como un Producto Mínimo Viable (MVP) robusto y escalable, lo que deja la puerta abierta a múltiples ampliaciones funcionales. A corto y medio plazo, se proponen las siguientes líneas de trabajo para evolucionar la plataforma:

\begin{itemize}
    \item \textbf{Personalización Avanzada de Escalas \hyperref[tab:esp_cu14]{(CU14)}:} Implementar la funcionalidad, ya diseñada e implementada a nivel de base de datos y \textit{backend} pero tuvo que ser pospuesta durante la implementación del frontend para poder priorizar otras funcionalidades principales, que permita al usuario reescribir los textos descriptivos de la escala del 1 al 10 y seleccionar conjuntos de iconos personalizados para representar cada estado de ánimo.
    \item \textbf{Nuevas Estadísticas:} En la versión actual de la aplicación, el usuario puede ponerle una calificación individual a cada etiqueta/actividad por separado (por ejemplo, puntuar un \texttt{2/10} en estudios). Aunque este dato ya se guarda correctamente en la base de datos, todavía no se utiliza para generar estadísticas. Una mejora importante para el futuro sería diseñar nuevas gráficas que utilicen estas notas para mostrar análisis mucho más profundos y detallados sobre el impacto real de cada actividad en el día a día.
    \item \textbf{Sistema de Notificaciones Activas:} Integrar un sistema de tareas programadas en el servidor que envíe recordatorios automáticos (vía correo electrónico) al dispositivo del usuario si, llegada una hora determinada, el sistema detecta que aún no ha registrado su jornada.
    \item \textbf{Aplicación Móvil Nativa:} Migrar el código de React a React Native para subir la app a la Google Play Store y Apple App Store.
    \item \textbf{Análisis de Sentimientos mediante Inteligencia Artificial:} Incorporar un modelo de Procesamiento de Lenguaje Natural en el \textit{backend} capaz de leer los campos del comentario libre, detectando palabras clave, sentimientos subyacentes o patrones de estrés que complementen el análisis numérico de las etiquetas.
\end{itemize}